\chapter{Orbital Perbatasan dan Kereaktifan}\label{ch:b2:frontier-orbitals}

Sebotol siklopentadiena yang dibiarkan di rak perlahan-lahan berubah
menjadi dimernya: dua molekul \omterm{def:b2:frontier-orbitals:diels-alder}{diena} yang sama bergabung menjadi senyawa
bisiklik berjembatan, dan hanya satu dari stereoisomer-stereoisomer yang
mungkin yang terbentuk. Tidak ada apa pun dalam panah-panah lengkung jilid
Tahun 1 yang meramalkan stereoisomer mana. Dua orbital meramalkannya:
orbital terisi tertinggi satu molekul dan orbital kosong terendah molekul
yang lain. Bab ini mengubah interaksi dua tingkat pada
\cref{ch:b2:diatomic-mos} menjadi alat untuk memahami kereaktifan --- atom
mana yang diserang nukleofil, dari sisi mana, dan mengapa \omterm{def:b2:frontier-orbitals:diels-alder}{diena} dan alkena
bergabung menjadi cincin dalam satu tahap.

\begin{recall}
\Cref{ch:b2:diatomic-mos}: dua orbital yang berinteraksi, tolakan empat
elektron. \Cref{ch:b2:huckel}: tingkat dan koefisien Hückel sistem
terkonjugasi. Jilid Tahun 1: nukleofil dan elektrofil, panah lengkung,
kendali kinetik, reaksi stereospesifik, stereoselektif, dan regioselektif,
adisi elektrofilik pada alkena, substitusi nukleofilik.
\end{recall}

\section{Dua orbital, sekali lagi}

Ketika dua molekul saling mendekat, orbital-orbitalnya bertumpang-tindih
dan berinteraksi. Sebagian besar interaksi itu lemah, karena
orbital-orbitalnya berjauhan energinya; bentuk perturbatif dari hasil dua
tingkat lalu sudah memadai.

\begin{theorem}[Interaksi lemah dua tingkat]\label{thm:b2:frontier-orbitals:perturbation}
Dua orbital berenergi $\varepsilon_1 < \varepsilon_2$, yang dikopel oleh \omterm{def:b2:diatomic-mos:integrals}{integral
resonansi} $\beta$ dengan $|\beta| \ll \Delta\varepsilon = \varepsilon_2 - \varepsilon_1$
(tumpang-tindih diabaikan), bergeser menjadi
\[
  \varepsilon_1 - \frac{\beta^2}{\Delta\varepsilon}, \qquad \varepsilon_2 + \frac{\beta^2}{\Delta\varepsilon}
\]
sampai orde pertama dalam $\beta^2/\Delta\varepsilon^2$: tingkat yang lebih
rendah turun dan yang lebih tinggi naik sebesar nilai yang sama, makin kecil
makin besar celahnya.
\end{theorem}

\begin{proof}
Menurut \cref{thm:b2:diatomic-mos:heteronuclear}, tingkat-tingkat eksaknya adalah $\bar\varepsilon \mp
\sqrt{\Delta\varepsilon^2/4 + \beta^2}$. Dengan $u = 4\beta^2/\Delta\varepsilon^2 \ll 1$, $\sqrt{\Delta\varepsilon^2/4
+ \beta^2} = \frac{\Delta\varepsilon}{2}\sqrt{1 + u} \approx \frac{\Delta\varepsilon}{2}(1 + u/2) =
\frac{\Delta\varepsilon}{2} + \frac{\beta^2}{\Delta\varepsilon}$, dan $\bar\varepsilon \mp \Delta\varepsilon/2 = \varepsilon_1,
\varepsilon_2$.
\end{proof}

\begin{proposition}[Dua elektron menarik, empat menolak]\label{prop:b2:frontier-orbitals:four-electron}
Jika kedua orbital yang berinteraksi seluruhnya memuat dua elektron, sistem
itu distabilkan sekitar $2\beta^2/\Delta\varepsilon$. Jika keduanya memuat
empat, sistem itu didestabilkan.
\end{proposition}

\begin{proof}
Dua elektron masuk ke tingkat yang diturunkan: keuntungannya
$2\beta^2/\Delta\varepsilon$. Dengan empat, tingkat atas juga terisi; sampai
orde \cref{thm:b2:frontier-orbitals:perturbation} kedua pergeseran saling
meniadakan, dan dengan tumpang-tindih yang diperhitungkan, tingkat atas
naik lebih jauh daripada turunnya tingkat bawah
(\cref{prop:b2:diatomic-mos:asymmetry}): sebuah tolakan bersih.
\end{proof}

\begin{omfigure}
\begin{tikzpicture}[font=\small]
  \begin{scope}
    \pic (a) at (0,0) {omlevel={ud}{0.7}};
    \pic (b) at (3,1.4) {omlevel={empty}{0.7}};
    \pic (lo) at (1.5,-0.5) {omlevel={ud}{0.7}};
    \pic (hi) at (1.5,1.9) {omlevel={empty}{0.7}};
    \draw[omcorr, densely dashed] (a-r) -- (lo-l) (a-r) -- (hi-l) (b-l) -- (lo-r) (b-l) -- (hi-r);
    \node[below=6pt] at (1.5,-0.6) {dua elektron: distabilkan};
    \node[left] at (a-l) {terisi};
    \node[right] at (b-r) {kosong};
  \end{scope}
  \begin{scope}[xshift=7cm]
    \pic (a) at (0,0) {omlevel={ud}{0.7}};
    \pic (b) at (3,0.6) {omlevel={ud}{0.7}};
    \pic (lo) at (1.5,-0.5) {omlevel={ud}{0.7}};
    \pic (hi) at (1.5,1.25) {omlevel={ud}{0.7}};
    \draw[omcorr, densely dashed] (a-r) -- (lo-l) (a-r) -- (hi-l) (b-l) -- (lo-r) (b-l) -- (hi-r);
    \node[below=6pt] at (1.5,-0.6) {empat elektron: tolak-menolak};
    \node[left] at (a-l) {terisi};
    \node[right] at (b-r) {terisi};
  \end{scope}
\end{tikzpicture}

{\small Interaksi dua orbital molekul-molekul yang bertetangga. Orbital
terisi yang berhadapan dengan orbital kosong memberikan stabilisasi
bersih; dua orbital terisi saling menolak, dengan tingkat atas naik lebih
jauh daripada turunnya tingkat bawah. Tolakan inilah asal-usul halangan
sterik antarmolekul.}
\end{omfigure}

\section{Orbital perbatasan}

\begin{definition}[Orbital perbatasan]\label{def:b2:frontier-orbitals:frontier}
\emph{HOMO}\index{HOMO} (\omterm{def:b2:diatomic-mos:lcao}{orbital molekul} terisi tertinggi) dan
\emph{LUMO}\index{LUMO} (\omterm{def:b2:diatomic-mos:lcao}{orbital molekul} kosong terendah) sebuah molekul
disebut \emph{orbital perbatasan}\index{orbital perbatasan} molekul itu.
\end{definition}

\begin{proposition}[Interaksi yang dominan]\label{prop:b2:frontier-orbitals:dominant}
Di antara dua molekul berkulit tertutup, interaksi penstabil yang paling
penting adalah interaksi antara \omterm{def:b2:frontier-orbitals:frontier}{HOMO} yang satu dan \omterm{def:b2:frontier-orbitals:frontier}{LUMO} yang lain; dari dua
pasangan semacam itu, pasangan dengan celah yang lebih kecil dominan.
Nukleofil bereaksi melalui \omterm{def:b2:frontier-orbitals:frontier}{HOMO-nya}, elektrofil melalui \omterm{def:b2:frontier-orbitals:frontier}{LUMO-nya}.
\end{proposition}

\begin{proof}[Argumen]
Pasangan terisi--terisi saling menolak
(\cref{prop:b2:frontier-orbitals:four-electron}) dan pasangan kosong--kosong
tidak memuat elektron: hanya pasangan terisi--kosong yang menstabilkan,
masing-masing sebesar $2\beta^2/\Delta\varepsilon$. \omterm{def:b2:frontier-orbitals:frontier}{HOMO} adalah tingkat terisi
tertinggi dan \omterm{def:b2:frontier-orbitals:frontier}{LUMO} tingkat kosong terendah, sehingga pasangan HOMO--LUMO
mempunyai celah terkecil di antara pasangan-pasangan itu; menurut teorema
itu, celah terkecil memberikan suku terbesar.
\end{proof}

\begin{definition}[Kendali muatan dan kendali orbital]\label{def:b2:frontier-orbitals:control}
Sebuah reaksi berada di bawah \emph{kendali muatan}\index{kendali muatan}
jika tarikan antara muatan-muatan kedua pasangan menentukan di mana
keduanya bereaksi, dan di bawah \emph{kendali orbital}\index{kendali orbital}
jika tumpang-tindih orbital-orbital perbatasan yang menentukannya, dengan
serangan yang terjadi di tempat koefisien-koefisiennya paling besar.
\end{definition}

\begin{definition}[Spesies keras dan lunak]\label{def:b2:frontier-orbitals:hard-soft}
Sebuah \emph{nukleofil keras}\index{nukleofil keras} atau \emph{elektrofil
keras}\index{elektrofil keras} berukuran kecil, membawa muatan yang
terpusat, dan mempunyai \omterm{def:b2:frontier-orbitals:frontier}{HOMO} (berturut-turut \omterm{def:b2:frontier-orbitals:frontier}{LUMO}) yang jauh dari orbital
pasangannya: spesies itu bereaksi di bawah \omterm{def:b2:frontier-orbitals:control}{kendali muatan}. Sebuah
\emph{nukleofil lunak}\index{nukleofil lunak} atau \emph{elektrofil
lunak}\index{elektrofil lunak} berukuran besar dan mudah terpolarisasi,
dengan \omterm{def:b2:frontier-orbitals:frontier}{HOMO} yang tinggi (berturut-turut \omterm{def:b2:frontier-orbitals:frontier}{LUMO} yang rendah): spesies itu
bereaksi di bawah \omterm{def:b2:frontier-orbitals:control}{kendali orbital}. Yang keras lebih suka bereaksi dengan
yang keras, yang lunak dengan yang lunak.
\end{definition}

Ion hidroksida dan ion fluorida adalah \omterm{def:b2:frontier-orbitals:hard-soft}{nukleofil keras}, iodida dan tiolat
lunak; proton dan karbokation adalah \omterm{def:b2:frontier-orbitals:hard-soft}{elektrofil keras}, karbon haloalkana
lunak.

\begin{definition}[Nukleofil ambiden]\label{def:b2:frontier-orbitals:ambident}
Sebuah \emph{nukleofil ambiden}\index{nukleofil ambiden} mempunyai dua atom
nukleofilik yang dihubungkan oleh konjugasi, yang masing-masing dapat
menyerang elektrofil.
\end{definition}

Ion sianida diserang pada nitrogen, tempat muatan negatifnya lebih besar,
oleh \omterm{def:b2:frontier-orbitals:hard-soft}{elektrofil keras}, dan pada karbon, tempat \omterm{def:b2:frontier-orbitals:frontier}{HOMO-nya} mempunyai koefisien
yang lebih besar, oleh \omterm{def:b2:frontier-orbitals:hard-soft}{elektrofil lunak}: haloalkana menghasilkan nitril
\ce{R-C#N}. Ion-ion enolat, yang dijumpai pada \cref{ch:b2:enolates-aldol},
bereaksi pada karbon atau pada oksigen dengan cara yang sama.

\begin{method}[Analisis orbital perbatasan]\label{met:b2:frontier-orbitals:fmo}
\begin{enumerate}
  \item Kenali nukleofilnya (\omterm{def:b2:frontier-orbitals:frontier}{HOMO} tinggi) dan elektrofilnya (\omterm{def:b2:frontier-orbitals:frontier}{LUMO} rendah).
  \item Gambarlah kedua \omterm{def:b2:frontier-orbitals:frontier}{orbital perbatasan} beserta koefisien dan tandanya.
  \item Reaksi terjadi di tempat tumpang-tindihnya paling besar: atom-atom
        dengan koefisien terbesar, cuping-cuping bertanda sama yang saling
        berhadapan.
  \item Arah pendekatan mengikuti bentuk \omterm{def:b2:frontier-orbitals:frontier}{LUMO} (tegak lurus bidang
        \ce{C=O}, di belakang ikatan \ce{C-X}).
  \item Celah HOMO--LUMO yang lebih kecil berarti reaksi yang lebih cepat
        (dengan hal-hal lain tetap sama).
\end{enumerate}
\end{method}

Arah serangan mengikuti \omterm{def:b2:frontier-orbitals:frontier}{LUMO}. \omterm{def:b2:frontier-orbitals:frontier}{LUMO} gugus karbonil adalah orbital
$\pi^*$-nya, yang lebih besar pada karbon dan tegak lurus bidang: nukleofil
mendekati karbon dari atas atau dari bawah bidang, bukan sepanjang sumbu
\ce{C=O}. \omterm{def:b2:frontier-orbitals:frontier}{LUMO} haloalkana adalah orbital $\sigma^*$ ikatan \ce{C-X}, yang
cuping besarnya pada karbon mengarah menjauhi X: nukleofil menyerang dari
belakang, dan karbonnya terinversi --- stereokimia substitusi bimolekular
dari jilid Tahun 1.

\section{Reaksi Diels--Alder}

\begin{definition}[Reaksi serempak]\label{def:b2:frontier-orbitals:concerted}
Yang disebut \emph{reaksi serempak}\index{reaksi serempak} adalah reaksi
yang membentuk dan memutus semua ikatannya dalam satu tahap, melalui satu
keadaan transisi, tanpa zat antara.
\end{definition}

\begin{definition}[Reaksi Diels--Alder]\label{def:b2:frontier-orbitals:diels-alder}
Yang disebut \emph{reaksi Diels--Alder}\index{reaksi Diels--Alder} adalah
adisi serempak sebuah \emph{diena}\index{diena} terkonjugasi, dalam
konformasi s-cis-nya, pada alkena atau alkuna, yaitu
\emph{dienofil}\index{dienofil}, yang membentuk cincin beranggota enam
dengan dua ikatan $\sigma$ baru dan satu ikatan $\pi$ baru.
\end{definition}

\begin{omfigure}
\begin{tikzpicture}[font=\small, >={Stealth}, scale=1.25]
  % diene s-cis: C1 (0,1.6) C2 (0.6,2.4) C3 (1.6,2.4) C4 (2.2,1.6); dienophile below
  \draw[thick] (0,1.6) -- (0.6,2.4) -- (1.6,2.4) -- (2.2,1.6);
  \draw[thick] (0.2,1.62) -- (0.65,2.22);
  \draw[thick] (2.0,1.62) -- (1.55,2.22);
  \draw[thick] (0.2,0) -- (2.0,0); \draw[thick] (0.35,0.14) -- (1.85,0.14);
  \node[font=\scriptsize, anchor=east] at (-0.05,1.6) {1};
  \node[font=\scriptsize, anchor=south east] at (0.5,2.42) {2};
  \node[font=\scriptsize, anchor=south] at (1.6,2.42) {3};
  \node[font=\scriptsize, anchor=west] at (2.25,1.6) {4};
  % arrows: dienophile pi -> C1..C(dienophile); C1=C2 pi -> C2-C3; C3=C4 pi -> C4-C
  \draw[->, thick, omProp] (1.1,0.22) .. controls (0.9,0.9) and (0.3,0.9) .. (0.12,1.42);
  \draw[->, thick, omProp] (0.45,1.95) .. controls (0.6,2.75) and (1.0,2.8) .. (1.1,2.48);
  \draw[->, thick, omProp] (1.78,1.92) .. controls (2.5,1.6) and (2.4,0.6) .. (2.05,0.2);
  \node at (3.3,1.2) {$\longrightarrow$};
  \begin{scope}[xshift=4.2cm]
    \draw[thick] (0,1.6) -- (0.6,2.4) -- (1.6,2.4) -- (2.2,1.6) -- (2.0,0) -- (0.2,0) -- cycle;
    \draw[thick] (0.72,2.26) -- (1.48,2.26);
  \end{scope}
\end{tikzpicture}

{\small \omterm{def:b2:frontier-orbitals:diels-alder}{Reaksi Diels--Alder} butadiena dan etena, yang dituliskan dengan
panah lengkung: tiga pasang elektron bergerak sekaligus, dalam satu tahap
siklik yang serempak; dua ikatan $\sigma$ terbentuk di ujung-ujung \omterm{def:b2:frontier-orbitals:diels-alder}{diena},
dan ikatan $\pi$ berpindah ke tengahnya.}
\end{omfigure}

Interaksi yang dominan adalah interaksi antara \omterm{def:b2:frontier-orbitals:frontier}{HOMO} \omterm{def:b2:frontier-orbitals:diels-alder}{diena}, $\psi_2$, dan \omterm{def:b2:frontier-orbitals:frontier}{LUMO}
\omterm{def:b2:frontier-orbitals:diels-alder}{dienofil}, $\pi^*$. Cuping-cuping keduanya di ujung-ujung \omterm{def:b2:frontier-orbitals:diels-alder}{diena} dan pada
kedua karbon \omterm{def:b2:frontier-orbitals:diels-alder}{dienofil} bertanda sesuai ketika \omterm{def:b2:frontier-orbitals:diels-alder}{dienofil} mendekati muka \omterm{def:b2:frontier-orbitals:diels-alder}{diena}:
kedua ikatan baru dapat terbentuk sekaligus, pada muka yang sama dari
masing-masing pasangan.

\begin{omfigure}
\begin{tikzpicture}[font=\small]
  \foreach \x/\c in {0/0.602, 1/0.372, 2/-0.372, 3/-0.602} {
    \pic at (\x,2) {omp={1.5*\c}{90}};
  }
  \draw[black!45] (-0.2,2) -- (3.2,2);
  \node[anchor=west] at (3.6,2) {HOMO diena ($\psi_2$)};
  \pic at (0,0) {omp={-0.85}{90}}; \pic at (3,0) {omp={0.85}{90}};
  \draw[black!45] (-0.2,0) -- (3.2,0);
  \node[anchor=west] at (3.6,0) {LUMO dienofil ($\pi^*$)};
  \draw[omProp, thick, densely dashed] (0,1.45) -- (0,0.6);
  \draw[omProp, thick, densely dashed] (3,1.45) -- (3,0.6);
  \node[font=\scriptsize, omProp, anchor=east] at (-0.15,1.0) {sefase};
  \node[font=\scriptsize, omProp, anchor=west] at (3.15,1.0) {sefase};
\end{tikzpicture}

{\small Orbital-orbital perbatasan butadiena dan etena, dengan cuping yang
sebanding dengan koefisien Hückel-nya. Ketika saling berhadapan,
cuping-cuping ujung \omterm{def:b2:frontier-orbitals:frontier}{HOMO} \omterm{def:b2:frontier-orbitals:diels-alder}{diena} dan cuping-cuping \omterm{def:b2:frontier-orbitals:frontier}{LUMO} \omterm{def:b2:frontier-orbitals:diels-alder}{dienofil} mempunyai
tanda yang sama di kedua ujungnya: kedua ikatan $\sigma$ dapat terbentuk
bersama-sama, pada muka yang sama dari masing-masing pasangan.}
\end{omfigure}

\begin{proposition}[Stereospesifisitas]\label{prop:b2:frontier-orbitals:da-stereospecific}
\omterm{def:b2:frontier-orbitals:diels-alder}{Reaksi Diels--Alder} bersifat stereospesifik: substituen-substituen yang cis
pada \omterm{def:b2:frontier-orbitals:diels-alder}{dienofil} tetap cis dalam produk, substituen-substituen trans tetap
trans.
\end{proposition}

\begin{proof}
Reaksinya serempak dan kedua ikatan baru terbentuk pada muka yang sama
dari \omterm{def:b2:frontier-orbitals:diels-alder}{dienofil}: kedua karbonnya tidak pernah berotasi satu terhadap yang
lain, dan posisi relatif substituen-substituennya terbawa ke dalam cincin.
\end{proof}

\begin{proposition}[Regioselektivitas]\label{prop:b2:frontier-orbitals:da-regio}
Dengan pasangan-pasangan yang tidak simetris, produk utamanya
menghubungkan atom \omterm{def:b2:frontier-orbitals:diels-alder}{diena} dengan koefisien terbesar dalam \omterm{def:b2:frontier-orbitals:frontier}{HOMO-nya} dengan
atom \omterm{def:b2:frontier-orbitals:diels-alder}{dienofil} dengan koefisien terbesar dalam \omterm{def:b2:frontier-orbitals:frontier}{LUMO-nya}.
\end{proposition}

\begin{proof}[Argumen]
Dalam keadaan transisi, kedua ikatan baru belum sama-sama terbentuk;
stabilisasi $2\beta^2/\Delta\varepsilon$ lebih besar jika koefisien-koefisien
yang lebih besar bertumpang-tindih ($\beta$ bertambah seiring hasil kali
koefisien atom-atom yang bertumpang-tindih), sehingga orientasi itu dicapai
pada energi yang lebih rendah.
\end{proof}

\begin{omfigure}
\begin{tikzpicture}
\begin{axis}[ybar, width=0.47\linewidth, height=4.6cm, ymin=-0.8, ymax=0.8,
  xtick={0,1,2,3,4}, xticklabels={O,C1,C2,C3,C4}, axis x line=bottom, axis y line=left,
  extra y ticks={0}, extra y tick style={grid=major, grid style={black!50}}, enlarge x limits=0.15,
  tick label style={font=\scriptsize}, ylabel={koefisien}, label style={font=\small},
  title={\small HOMO 1-metoksibutadiena}, bar width=8pt]
\addplot[fill=omDef!50, draw=omDef] table[x=atom, y=c] {figdata/out/bachelor-2-huckel-c-diene.dat};
\end{axis}
\end{tikzpicture}\hfill
\begin{tikzpicture}
\begin{axis}[ybar, width=0.47\linewidth, height=4.6cm, ymin=-0.8, ymax=0.8,
  xtick={1,2,3,4}, xticklabels={C$\beta$,C$\alpha$,C(O),O}, axis x line=bottom, axis y line=left,
  extra y ticks={0}, extra y tick style={grid=major, grid style={black!50}}, enlarge x limits=0.15,
  tick label style={font=\scriptsize}, ylabel={koefisien}, label style={font=\small},
  title={\small LUMO propenal}, bar width=8pt]
\addplot[fill=omThm!45, draw=omThm] table[x=atom, y=c] {figdata/out/bachelor-2-huckel-c-dienophile.dat};
\end{axis}
\end{tikzpicture}

{\small Model Hückel dengan parameter heteroatom ($\alpha_{\ce{O}} = \alpha + 2\beta$,
$\beta_{\ce{CO}} = 0.8\beta$ untuk oksigen eter; $\alpha_{\ce{O}} = \alpha + \beta$, $\beta_{\ce{CO}} =
\beta$ untuk oksigen karbonil; sebuah model). \omterm{def:b2:frontier-orbitals:frontier}{HOMO} \omterm{def:b2:frontier-orbitals:diels-alder}{diena} paling besar pada
C4, ujung yang jauh dari gugus metoksi (0.60 dibandingkan dengan 0.51);
\omterm{def:b2:frontier-orbitals:frontier}{LUMO} \omterm{def:b2:frontier-orbitals:diels-alder}{dienofil} paling besar pada karbon $\beta$ (0.66). Keduanya berikatan
bersama: gugus metoksi dan gugus aldehida berakhir pada karbon-karbon
cincin yang bertetangga.}
\end{omfigure}

Gugus donor menaikkan \omterm{def:b2:frontier-orbitals:frontier}{HOMO} \omterm{def:b2:frontier-orbitals:diels-alder}{diena} ($m = 0.48$ alih-alih 0.62) dan gugus
akseptor menurunkan \omterm{def:b2:frontier-orbitals:frontier}{LUMO} \omterm{def:b2:frontier-orbitals:diels-alder}{dienofil} ($m = -0.35$ alih-alih $-1$): celahnya
menyempit, dan pasangan semacam itu bereaksi jauh lebih cepat daripada
butadiena dengan etena. \omterm{def:b2:frontier-orbitals:diels-alder}{Diena} yang kaya elektron dan \omterm{def:b2:frontier-orbitals:diels-alder}{dienofil} yang miskin
elektron adalah kombinasi yang khas.

\begin{definition}[Adduk endo dan exo]\label{def:b2:frontier-orbitals:endo}
Jika sebuah \omterm{def:b2:frontier-orbitals:diels-alder}{diena} siklik beradisi pada \omterm{def:b2:frontier-orbitals:diels-alder}{dienofil} yang membawa substituen
tak jenuh, \emph{adduk endo}\index{adduk endo} mempunyai substituen itu
mengarah ke ikatan rangkap yang baru terbentuk (di bawah \omterm{def:b2:frontier-orbitals:diels-alder}{diena} dalam
keadaan transisi), sedangkan \emph{adduk exo}\index{adduk exo} mengarah
menjauhinya. Yang disebut \emph{aturan endo}\index{aturan endo} menyatakan
bahwa adduk endo terbentuk lebih cepat.
\end{definition}

\begin{proposition}[Aturan endo]\label{prop:b2:frontier-orbitals:endo}
Di bawah kendali kinetik, \omterm{def:b2:frontier-orbitals:endo}{adduk endo} adalah produk utama, walaupun \omterm{def:b2:frontier-orbitals:endo}{adduk
exo} sering lebih stabil.
\end{proposition}

\admitted

Penjelasan yang lazim adalah tumpang-tindih sekunder, pada pendekatan
endo, antara cuping-cuping sistem $\pi$ substituen dan cuping-cuping
atom-atom tengah \omterm{def:b2:frontier-orbitals:frontier}{HOMO} \omterm{def:b2:frontier-orbitals:diels-alder}{diena}; tumpang-tindih itu menurunkan keadaan transisi
endo tanpa membentuk ikatan. Aturan ini empiris dan mempunyai
pengecualian; jilid Tahun 3 membahas reaksi-reaksi ini dengan simetri
orbital-orbitalnya.

\begin{omfigure}
\begin{tikzpicture}[font=\small]
  \foreach \xs/\mode in {0/endo, 5/exo} {
    \begin{scope}[xshift=\xs cm]
      % cyclopentadiene seen from the side: C1, C2=C3, C4, CH2 bridge on top
      \draw[thick] (0,2) -- (0.6,2.5) -- (1.8,2.5) -- (2.4,2);
      \draw[thick] (0.75,2.38) -- (1.65,2.38);
      \draw[thick] (0,2) -- (1.2,3.2) -- (2.4,2);
      \node[font=\scriptsize, anchor=south] at (1.2,3.2) {\ce{CH2}};
      % dienophile
      \draw[thick] (0.2,0) -- (2.2,0); \draw[thick] (0.35,0.1) -- (2.05,0.1);
      \draw[omProp, densely dashed, thick] (0,1.9) -- (0.2,0.15);
      \draw[omProp, densely dashed, thick] (2.4,1.9) -- (2.2,0.15);
      \node[anchor=north] at (1.2,-1.3) {\mode};
    \end{scope}
  }
  % endo: carbonyls tucked under the diene
  \draw[thick] (0.2,0) -- (0.75,0.95) (2.2,0) -- (1.65,0.95);
  \node[font=\scriptsize, fill=white, inner sep=1pt] at (0.75,1.1) {C=O};
  \node[font=\scriptsize, fill=white, inner sep=1pt] at (1.65,1.1) {C=O};
  \draw[omThm, densely dotted, thick] (0.75,1.25) -- (0.65,2.35);
  \draw[omThm, densely dotted, thick] (1.65,1.25) -- (1.75,2.35);
  % exo: carbonyls pointing away
  \draw[thick] (5.2,0) -- (4.95,-0.85) (7.2,0) -- (7.45,-0.85);
  \node[font=\scriptsize] at (4.95,-1.0) {C=O};
  \node[font=\scriptsize] at (7.45,-1.0) {C=O};
\end{tikzpicture}

{\small Pendekatan endo dan exo anhidrida maleat pada siklopentadiena,
skematis. Ikatan-ikatan yang sedang terbentuk digambar putus-putus. Pada
pendekatan endo, gugus-gugus karbonil terletak di bawah \omterm{def:b2:frontier-orbitals:diels-alder}{diena}, dan orbital
$\pi$-nya bertumpang-tindih dengan karbon-karbon tengah \omterm{def:b2:frontier-orbitals:diels-alder}{diena} (titik-titik,
tumpang-tindih sekunder); pada pendekatan exo, gugus-gugus itu mengarah
menjauh.}
\end{omfigure}

\begin{method}[Menggambar produk Diels--Alder]\label{met:b2:frontier-orbitals:diels-alder}
\begin{enumerate}
  \item Letakkan \omterm{def:b2:frontier-orbitals:diels-alder}{diena} dalam konformasi s-cis-nya, dengan \omterm{def:b2:frontier-orbitals:diels-alder}{dienofil}
        menghadap ujung-ujungnya.
  \item Gambarlah kedua ikatan $\sigma$ baru dari ujung-ujung \omterm{def:b2:frontier-orbitals:diels-alder}{diena} ke
        karbon-karbon \omterm{def:b2:frontier-orbitals:diels-alder}{dienofil}, dan ikatan $\pi$ baru antara C2 dan C3 \omterm{def:b2:frontier-orbitals:diels-alder}{diena}.
  \item Pertahankan substituen-substituen \omterm{def:b2:frontier-orbitals:diels-alder}{dienofil} pada sisi yang sama (cis
        tetap cis).
  \item Untuk pasangan yang tidak simetris, hubungkan koefisien-koefisien
        terbesar (produk ``orto'' atau ``para'' untuk \omterm{def:b2:frontier-orbitals:diels-alder}{diena} tersubstitusi
        pada posisi 1 atau 2).
  \item Dengan \omterm{def:b2:frontier-orbitals:diels-alder}{diena} siklik, letakkan substituen tak jenuh pada posisi
        endo.
\end{enumerate}
\end{method}

\begin{inthelab}[Perengkahan disiklopentadiena]
Siklopentadiena dijual dalam bentuk dimernya dan diregenerasi tepat sebelum
dipakai: dimernya dipanaskan sampai titik didihnya (sekitar
\qty{170}{\degreeCelsius}) dalam labu yang dilengkapi kolom fraksionasi,
tempat dimer itu terpecah kembali menjadi monomer (reaksi retro-Diels--Alder),
yang terdistilasi pada \qty{41}{\degreeCelsius} dan ditampung dalam
penampung yang didinginkan dengan es. Monomer itu kembali berdimerisasi
dalam beberapa jam pada suhu kamar dan langsung dipakai. Kedua senyawa itu
mudah terbakar dan berbahaya; pekerjaan dilakukan di dalam lemari asam.
\end{inthelab}

\begin{safety}{\ghs[5mm]{GHS02}\,\ghs[5mm]{GHS06}\,\ghs[5mm]{GHS08}}
Siklopentadiena dan dimernya sangat mudah terbakar dan beracun jika
terhirup; anhidrida maleat bersifat korosif dan menyebabkan sensitisasi
pernapasan. Senyawa-senyawa itu ditangani di dalam lemari asam, dengan
sarung tangan dan pelindung mata, jauh dari api.
% ledger: ghs:cyclopentadiene, ghs:dicyclopentadiene, ghs:maleic-anhydride, bp:cyclopentadiene, bp:dicyclopentadiene
\end{safety}

\section{Latihan}

\begin{exercise}[$\star$]\label{exo:b2:frontier-orbitals:1}
Berikan \omterm{def:b2:frontier-orbitals:frontier}{HOMO} dan \omterm{def:b2:frontier-orbitals:frontier}{LUMO} (energi dalam satuan Hückel) etena, anion alil, dan
butadiena.
\end{exercise}

\begin{exercise}[$\star$]\label{exo:b2:frontier-orbitals:2}
Golongkan sebagai keras atau lunak: \ce{OH-}, \ce{I-}, \ce{H+}, tiolat
\ce{RS-}, karbon iodometana, \ce{F-}.
\end{exercise}

\begin{exercise}[$\star$]\label{exo:b2:frontier-orbitals:3}
\omterm{def:b2:frontier-orbitals:diels-alder}{Diena} manakah di antara berikut ini yang dapat mengambil bagian dalam
\omterm{def:b2:frontier-orbitals:diels-alder}{reaksi Diels--Alder}: buta-1,3-diena, siklopentadiena, penta-1,4-diena,
(2\textit{Z},4\textit{Z})-heksa-2,4-diena dalam bentuk s-cis-nya?
\end{exercise}

\begin{exercise}[$\star$]\label{exo:b2:frontier-orbitals:4}
Gambarlah produk butadiena dengan propenanitril \ce{CH2=CH-C#N}.
\end{exercise}

\begin{exercise}[$\star\star$]\label{exo:b2:frontier-orbitals:5}
Dua orbital pada $-10$ dan \qty{-2}{eV} berinteraksi dengan
$\beta = \qty{-1.0}{eV}$ (data latihan). Hitunglah stabilisasi perturbatif
dua elektron, lalu stabilisasi eksaknya.
\end{exercise}

\begin{exercise}[$\star\star$]\label{exo:b2:frontier-orbitals:6}
Ion sianida bereaksi dengan iodometana menghasilkan etananitril
\ce{CH3-C#N}, tetapi dengan karbokation dalam pelarut yang sangat pengion
menghasilkan sebagian isonitril \ce{R-N#C}. Jelaskan.
\end{exercise}

\begin{exercise}[$\star\star$]\label{exo:b2:frontier-orbitals:7}
Dengan koefisien-koefisien dalam bab ini, ramalkan produk utama
1-metoksibutadiena dengan propenal.
\end{exercise}

\begin{exercise}[$\star\star$]\label{exo:b2:frontier-orbitals:8}
Gambarlah \omterm{def:b2:frontier-orbitals:endo}{adduk endo} siklopentadiena dengan metil propenoat.
\end{exercise}

\begin{exercise}[$\star\star$]\label{exo:b2:frontier-orbitals:9}
Mengapa anhidrida maleat bereaksi dengan siklopentadiena pada suhu kamar,
sedangkan etena memerlukan pemanasan di bawah tekanan?
\end{exercise}

\begin{exercise}[$\star\star\star$]\label{exo:b2:frontier-orbitals:10}
Siklopentadiena bereaksi dengan dimetil maleat (diester cis) dan dengan
dimetil fumarat (diester trans). Gambarlah kedua produknya dan nyatakan
berapa stereoisomer yang dihasilkan masing-masing.
\end{exercise}

\begin{exercise}[$\star\star\star$]\label{exo:b2:frontier-orbitals:11}
\omterm{def:b2:frontier-orbitals:diels-alder}{Reaksi Diels--Alder} mempunyai $\Delta_r H^\circ < 0$ dan $\Delta_r S^\circ < 0$.
Jelaskan tanda-tanda itu dan mengapa reaksi baliknya menjadi
menguntungkan pada suhu tinggi.
\end{exercise}

\begin{exercise}[$\star\star\star$]\label{exo:b2:frontier-orbitals:12}
Dalam hidrazina \ce{H2N-NH2}, setiap nitrogen membawa sepasang elektron
bebas. Dengan interaksi empat elektron, jelaskan mengapa molekul itu
menghindari konformasi yang kedua pasangan elektron bebasnya sejajar.
\end{exercise}

\section{Soal: Perengkahan Siklopentadiena}

\begin{problem}[{Soal akhir pekan --- dimerisasi siklopentadiena sebagai
reaksi Diels--Alder, termodinamika perengkahan, orbital perbatasan
siklopentadiena dan anhidrida maleat, serta stereokimia aduknya}]
\label{pb:b2:frontier-orbitals:1}
Data: titik didih, siklopentadiena \qty{41}{\degreeCelsius},
disiklopentadiena sekitar \qty{170}{\degreeCelsius}. Energi model Hückel
(sebuah model): siklopentadiena, yang diperlakukan sebagai satuan
butadiena, \omterm{def:b2:frontier-orbitals:frontier}{HOMO} $\alpha + 0.62\beta$ dan \omterm{def:b2:frontier-orbitals:frontier}{LUMO} $\alpha - 0.62\beta$; anhidrida
maleat, \omterm{def:b2:frontier-orbitals:frontier}{LUMO} sekitar $\alpha - 0.35\beta$.
% ledger: bp:cyclopentadiene, bp:dicyclopentadiene

\medskip
\noindent\textbf{Bagian I --- Dimer.}
\begin{enumerate}
  \item Gambarlah siklopentadiena dan tunjukkan bahwa satuan \omterm{def:b2:frontier-orbitals:diels-alder}{dienanya}
        tertahan dalam bentuk s-cis.
  \item Dalam dimerisasi, satu molekul berperan sebagai \omterm{def:b2:frontier-orbitals:diels-alder}{diena}. Apa peran
        molekul yang lain?
  \item Gambarlah dimernya dan beri nomor ikatan-ikatan barunya.
  \item Adduk manakah, endo atau exo, yang terbentuk pada suhu kamar, dan
        mengapa?
  \item Apakah dimerisasi itu stereospesifik? Jelaskan dengan mekanismenya.
  \item Mengapa dimerisasi itu tidak memerlukan katalis?
\end{enumerate}

\medskip
\noindent\textbf{Bagian II --- Termodinamika perengkahan.}
\begin{enumerate}[resume]
  \item Berikan tanda $\Delta_r H^\circ$ dimerisasi, dari ikatan-ikatan yang
        terbentuk dan yang diputus.
  \item Berikan tanda $\Delta_r S^\circ$.
  \item Tuliskan $\Delta_r G^\circ(T)$ dan tunjukkan bahwa tandanya berganti
        pada suhu $T_i$.
  \item Mengapa dimer stabil pada suhu kamar, dan monomer lebih disukai
        dalam keadaan panas?
  \item Dalam susunan alat perengkahan, mengapa mendistilasi monomer segera
        setelah terbentuk mendorong reaksi?
  \item Mengapa monomer harus dijaga tetap dingin dan segera dipakai?
\end{enumerate}

\medskip
\noindent\textbf{Bagian III --- \omterm{def:b2:frontier-orbitals:frontier}{Orbital perbatasan}.}
\begin{enumerate}[resume]
  \item Untuk siklopentadiena yang bereaksi dengan dirinya sendiri,
        hitunglah celah HOMO--LUMO dalam satuan $|\beta|$.
  \item Untuk siklopentadiena dengan anhidrida maleat, hitunglah celah
        antara \omterm{def:b2:frontier-orbitals:frontier}{HOMO} \omterm{def:b2:frontier-orbitals:diels-alder}{diena} dan \omterm{def:b2:frontier-orbitals:frontier}{LUMO} anhidrida.
  \item Pasangan manakah yang bereaksi lebih cepat? Jelaskan.
  \item Mengapa kedua gugus karbonil anhidrida maleat menurunkan \omterm{def:b2:frontier-orbitals:frontier}{LUMO-nya}?
  \item Orbital anhidrida maleat manakah yang memberikan tumpang-tindih
        sekunder pada pendekatan endo?
\end{enumerate}

\medskip
\noindent\textbf{Bagian IV --- Adduk dengan anhidrida maleat.}
\begin{enumerate}[resume]
  \item Gambarlah \omterm{def:b2:frontier-orbitals:endo}{adduk endo}.
  \item Tandailah pusat-pusat stereogeniknya.
  \item Apakah adduk itu kiral? (Carilah bidang cermin.)
  \item Mengapa kedua hidrogen bekas \omterm{def:b2:frontier-orbitals:diels-alder}{dienofil} berakhir pada sisi cincin yang
        sama?
  \item Seperti apa \omterm{def:b2:frontier-orbitals:endo}{adduk exo}?
  \item Dapatkah \omterm{def:b2:frontier-orbitals:endo}{adduk endo} dan exo saling berubah pada suhu kamar? Mengapa?
  \item Bagaimana anhidrida itu akan bereaksi dengan air, dan produk apa
        yang akan terbentuk?
  \item Nyatakan jumlah pusat stereogenik \omterm{def:b2:frontier-orbitals:endo}{adduk endo} siklopentadiena dan
        anhidrida maleat.
\end{enumerate}
\end{problem}
